\section*{Bloque A. Fundamentos teóricos y conceptuales}

\noindent\textit{Preguntas 1 a 30. Evalúan la comprensión analítica de los axiomas de las preferencias, la optimización del consumidor, la dualidad, la ecuación de Slutsky, las elasticidades y los conceptos básicos de bienestar.}

\q{La relación marginal de sustitución (RMS) de $x$ por $y$ se define como:}
\begin{opts}
    \item la pendiente absoluta de la recta presupuestaria en equilibrio.
    \item la relación de precios entre el bien $x$ y el bien $y$.
    \item unidades de $y$ que el consumidor cede por una de $x$ sin perder utilidad.
    \item la variación del ingreso real ante un aumento proporcional de precios.
\end{opts}

\q{Para la función de utilidad $U(x,y)=x^2y^3$, la RMS de $x$ por $y$ es:}
\begin{opts}
    \item $\dfrac{2y}{3x}$
    \item $\dfrac{3y}{2x}$
    \item $\dfrac{2x}{3y}$
    \item $\dfrac{3x}{2y}$
\end{opts}

\q{Una transformación monótona creciente de una función de utilidad:}
\begin{opts}
    \item altera las curvas de indiferencia y cambia toda la RMS.
    \item conserva intactas las utilidades marginales de cada bien.
    \item solo resulta válida para funciones de tipo Cobb-Douglas.
    \item conserva las curvas de indiferencia y la RMS, aunque cambie la $UMg$.
\end{opts}

\q{Si las preferencias son convexas, a lo largo de una curva de indiferencia la RMS:}
\begin{opts}
    \item es creciente y las curvas son cóncavas al origen.
    \item es decreciente y las curvas son convexas al origen.
    \item es constante a lo largo de toda la curva de indiferencia.
    \item tiende a infinito en cualquier punto de consumo interior.
\end{opts}

\q{Un aumento del precio de $x$, manteniendo constantes el ingreso y el precio de $y$, produce:}
\begin{opts}
    \item un desplazamiento paralelo hacia afuera de la recta presupuestaria.
    \item un giro de la recta presupuestaria hacia afuera sobre el eje $x$.
    \item un giro de la recta presupuestaria hacia adentro sobre el eje $x$.
    \item un desplazamiento paralelo hacia adentro.
\end{opts}

\q{Un impuesto específico de $t$ soles por unidad sobre el bien $x$ se representa mediante:}
\begin{opts}
    \item $p_xx+p_yy=I-t$.
    \item $(p_x+t)x+p_yy=I$.
    \item $p_xx+(p_y+t)y=I$.
    \item $p_xx+p_yy=I+tx$.
\end{opts}

\q{Un subsidio en especie que entrega $G$ unidades del bien $x$ difiere de un subsidio en efectivo de igual valor porque:}
\begin{opts}
    \item siempre genera mayor bienestar que un subsidio en efectivo.
    \item equivale analíticamente a duplicar el precio del bien $x$.
    \item deja inalterada la frontera del conjunto presupuestario.
    \item exige consumir al menos $G$ unidades de $x$ y puede restringir la elección.
\end{opts}

\q{En un óptimo interior con preferencias regulares, la condición de primer orden exige:}
\begin{opts}
    \item $RMS_{xy}=p_x/p_y$ junto con la condición de segundo orden.
    \item $RMS_{xy}>p_x/p_y$ en cualquier punto interior factible.
    \item $UMg_x=UMg_y$ con total independencia de precios relativos.
    \item el gasto monetario resulta idéntico en ambos bienes.
\end{opts}

\q{Si en una canasta interior se cumple $UMg_x/p_x>UMg_y/p_y$, el consumidor:}
\begin{opts}
    \item ya maximiza su utilidad en esa canasta interior factible.
    \item debe reasignar gasto desde el bien $x$ hacia el bien $y$.
    \item puede elevar su utilidad reasignando gasto de $y$ hacia $x$.
    \item se halla necesariamente en una solución de esquina óptima.
\end{opts}

\q{Para $U(x,y)=ax+by$, con $a,b>0$, si $p_x/p_y<a/b$, el consumidor elige:}
\begin{opts}
    \item solo $y$.
    \item solo $x$.
    \item una combinación interior con $x=y$.
    \item no consumir.
\end{opts}

\q{Para bienes perfectamente complementarios $U=\min(ax,by)$, las curvas de indiferencia tienen forma de:}
\begin{opts}
    \item líneas rectas con pendiente negativa.
    \item hipérbolas convexas.
    \item circunferencias.
    \item ángulos rectos (forma de L).
\end{opts}

\q{Con $U=\min(x,y)$, un aumento del precio de $x$:}
\begin{opts}
    \item reduce $x$ e $y$ en la misma proporción.
    \item no altera las cantidades demandadas.
    \item aumenta la cantidad demandada de $y$.
    \item deja la utilidad constante.
\end{opts}

\q{La función de utilidad cuasilineal $U=v(x)+y$ se caracteriza porque:}
\begin{opts}
    \item todos los bienes presentan elasticidad ingreso igual a uno.
    \item la demanda de $x$ no depende del ingreso cuando $y>0$.
    \item la relación marginal de sustitución es siempre constante.
    \item las curvas de indiferencia son líneas rectas paralelas.
\end{opts}

\q{Si las preferencias son homotéticas, las curvas de Engel:}
\begin{opts}
    \item tienen pendiente negativa en todo el rango de ingresos.
    \item son verticales e independientes del nivel de ingreso.
    \item son rayos desde el origen con participaciones de gasto fijas.
    \item solo existen para bienes inferiores de consumo básico.
\end{opts}

\q{Un bien inferior se define porque su elasticidad ingreso es:}
\begin{opts}
    \item mayor que 1.
    \item entre 0 y 1.
    \item exactamente 1.
    \item menor que 0.
\end{opts}

\q{Un bien de lujo se caracteriza porque su elasticidad ingreso es:}
\begin{opts}
    \item mayor que 1.
    \item entre 0 y 1.
    \item menor que 0.
    \item exactamente 0.
\end{opts}

\q{La demanda marshalliana de un bien se obtiene al:}
\begin{opts}
    \item minimizar el gasto sujeto a un nivel de utilidad dado.
    \item maximizar el ingreso disponible ante precios constantes.
    \item maximizar la utilidad sujeta a la restricción presupuestaria.
    \item igualar el precio de mercado al costo marginal unitario.
\end{opts}

\q{La demanda hicksiana o compensada se obtiene al:}
\begin{opts}
    \item maximizar la utilidad sujeto a un nivel de ingreso dado.
    \item minimizar el gasto sujeto a un nivel de utilidad dado.
    \item maximizar el excedente monetario neto del consumidor.
    \item derivar directamente la ecuación general de Slutsky.
\end{opts}

\q{La función de gasto $e(p,u)$ representa:}
\begin{opts}
    \item el gasto máximo factible para una canasta de precios dados.
    \item la función indirecta de utilidad evaluada en equilibrio.
    \item la curva ordinaria o demanda marshalliana del bien.
    \item el gasto mínimo para lograr utilidad $u$ a precios $p$.
\end{opts}

\q{La identidad de Roy permite obtener:}
\begin{opts}
    \item la demanda marshalliana a partir de la utilidad indirecta.
    \item la demanda hicksiana a partir de la función de gasto.
    \item la función de producción a partir del costo total mínimo.
    \item la elasticidad precio cruzada desde la elasticidad ingreso.
\end{opts}

\q{El lema de Shephard establece que:}
\begin{opts}
    \item $\partial v(p,I)/\partial p_i=x_i(p,I)$.
    \item $\partial e(p,u)/\partial u=v(p,I)$.
    \item $\partial e(p,u)/\partial p_i=h_i(p,u)$.
    \item $\partial U/\partial x_i=p_i$.
\end{opts}

\q{La ecuación de Slutsky descompone el efecto total de un cambio de precio en:}
\begin{opts}
    \item efecto escala y efecto técnico.
    \item efecto sustitución y efecto ingreso.
    \item efecto directo y efecto cruzado.
    \item efecto precio y efecto cantidad.
\end{opts}

\q{Para un bien normal, la curva de demanda compensada es:}
\begin{opts}
    \item más plana que la curva de demanda marshalliana ordinaria.
    \item idéntica en todo punto a la curva de demanda marshalliana.
    \item completamente vertical en todo el rango de precios.
    \item más empinada que la demanda marshalliana ordinaria.
\end{opts}

\q{Un bien Giffen es aquel que:}
\begin{opts}
    \item es inferior y el efecto ingreso domina, subiendo su demanda al subir $p$.
    \item es normal y su cantidad demandada sube al elevarse su propio precio.
    \item presenta demanda perfectamente elástica a lo largo de la curva.
    \item es un bien superior con elasticidad ingreso mayor que la unidad.
\end{opts}

\q{La elasticidad precio de una demanda vertical es:}
\begin{opts}
    \item $-1$.
    \item infinita.
    \item cero.
    \item positiva.
\end{opts}

\q{Si $|\varepsilon|>1$, un aumento del precio:}
\begin{opts}
    \item aumenta el ingreso total.
    \item reduce el ingreso total.
    \item no afecta el ingreso total.
    \item aumenta la cantidad demandada.
\end{opts}

\q{La homogeneidad de grado cero de la demanda marshalliana en $(p,I)$ implica que:}
\begin{opts}
    \item al duplicar el ingreso nominal, toda demanda se duplica.
    \item la curva de demanda ordinaria es elástica en todo punto.
    \item el efecto sustitución neto resulta idénticamente nulo.
    \item al duplicar precios e ingreso, la demanda no cambia.
\end{opts}

\q{La ley de Walras, en una economía de intercambio, establece que:}
\begin{opts}
    \item la suma del valor de excesos de demanda es cero en todos los mercados.
    \item cada mercado individual alcanza el equilibrio de forma aislada.
    \item los precios monetarios permanecen rígidos en el corto plazo.
    \item los consumidores agotan el ahorro antes de demandar bienes.
\end{opts}

\q{Las preferencias lexicográficas:}
\begin{opts}
    \item son continuas y admiten siempre una función de utilidad suave.
    \item son estrictamente convexas, monótonas y homotéticas.
    \item no son continuas ni admiten una función de utilidad.
    \item resultan equivalentes a preferencias de tipo Cobb-Douglas.
\end{opts}

\q{El axioma débil de preferencia revelada (WARP) sostiene que:}
\begin{opts}
    \item las preferencias del agente deben ser siempre transitivas.
    \item si $x$ se revela preferida a $y$, $y$ no puede revelarse preferida a $x$.
    \item el consumidor elige invariablemente la canasta más barata.
    \item la función de utilidad posee un significado puramente cardinal.
\end{opts}
